\chapter{Binary Shannon Balance and Information Geometry}
\label{ch:shannon}

\section{Entropy in natural units}
For a Bernoulli distribution with probabilities $\sigma$ and $1-\sigma$, the Shannon entropy in nats is
\begin{equation}
\Hbin(\sigma)
=-\sigma\ln\sigma-(1-\sigma)\ln(1-\sigma),
\qquad 0\leq\sigma\leq1.
\label{eq:binary-entropy}
\end{equation}
The endpoint convention $0\ln0=0$ is understood. Shannon's original axiomatic derivation identifies the logarithm as the additive measure compatible with composition of choices; the base fixes the unit \cite{Shannon1948}. In base $e$, a balanced binary choice has entropy $\ln2$.

\begin{theorem}[Unique binary entropy maximum]
The function $\Hbin$ is strictly concave on $(0,1)$ and has the unique maximum
\begin{equation}
\Hbin\!\left(\frac12\right)=\ln2.
\end{equation}
\end{theorem}

\begin{proof}
Differentiation gives
\begin{equation}
\Hbin'(\sigma)=\ln\!\left(\frac{1-\sigma}{\sigma}\right),
\end{equation}
so the only stationary point satisfies $1-\sigma=\sigma$, hence $\sigma=1/2$. A second differentiation yields
\begin{equation}
\Hbin''(\sigma)
=-\frac1\sigma-\frac1{1-\sigma}
=-\frac{1}{\sigma(1-\sigma)}<0.
\end{equation}
Therefore the stationary point is the unique global maximum. Substitution gives
\begin{equation}
\Hbin(1/2)
=-2\left(\frac12\ln\frac12\right)=\ln2.
\end{equation}
\end{proof}

The number $\ln2$ is therefore not being inserted by hand at the balanced point. It is the exact natural-unit information of one maximally uncertain binary distinction.

\section{Complement symmetry}
Define the complement map
\begin{equation}
C(\sigma)=1-\sigma.
\end{equation}
Then
\begin{equation}
\Hbin(C(\sigma))=\Hbin(\sigma),
\end{equation}
and the unique fixed point of $C$ is $1/2$. This is formally parallel to the zeta complement $s\mapsto1-s$. The parallel becomes exact only after the model identification
\begin{assumption}[Real-part/probability identification]
\label{ass:sigma}
Within the information-spinor representation, identify the Bernoulli coordinate with the real part of the zeta argument:
\begin{equation}
\sigma=\Ree s.
\end{equation}
\end{assumption}

This assumption is natural because both coordinates lie in $(0,1)$ in their respective critical domains and both are acted on by the same complement $\sigma\mapsto1-\sigma$. Nevertheless, standard analytic number theory does not declare $\Ree s$ to be a probability. The identification is a modelling step and remains subject to justification by an operator or integral representation.

\section{Fisher metric of the Bernoulli family}
The Bernoulli family has Fisher information
\begin{equation}
g(\sigma)=\frac{1}{\sigma(1-\sigma)}.
\end{equation}
Thus the infinitesimal Fisher--Rao line element is
\begin{equation}
\dd \ell^2=\frac{\dd\sigma^2}{\sigma(1-\sigma)}.
\label{eq:fisher-line}
\end{equation}
The metric diverges at the deterministic boundaries and is symmetric about $1/2$. Introduce an angular coordinate $\theta$ by
\begin{equation}
\sigma=\cos^2\!\left(\frac{\theta}{2}\right)
=\frac{1+\cos\theta}{2}.
\label{eq:sigma-theta}
\end{equation}
Then
\begin{equation}
\dd\sigma=-\frac12\sin\theta\,\dd\theta,
\qquad
\sigma(1-\sigma)=\frac14\sin^2\theta,
\end{equation}
so \eqref{eq:fisher-line} reduces to
\begin{equation}
\dd \ell^2=\dd\theta^2.
\end{equation}
The Bernoulli manifold is therefore an interval with a natural angular coordinate. The balanced point $\sigma=1/2$ corresponds to $\theta=\pi/2$.

This is already the polar angle of the Bloch-sphere representation used in the next chapter. The information metric and the quantum-state geometry are thus joined by the square-root embedding rather than by metaphor.

\section{Square-root amplitude embedding}
Define
\begin{equation}
\boldsymbol{a}(\sigma)=
\begin{pmatrix}
\sqrt{\sigma}\\
\sqrt{1-\sigma}
\end{pmatrix}.
\end{equation}
Then $\|\boldsymbol{a}\|^2=1$, so the probability simplex is embedded into the unit circle in amplitude coordinates. The complement exchanges the two components. The balanced point becomes
\begin{equation}
\boldsymbol{a}(1/2)=\frac{1}{\sqrt2}
\begin{pmatrix}1\\1\end{pmatrix}.
\end{equation}

The square-root map is important because coherent cancellation acts on amplitudes, not directly on probabilities. A phase can rotate one amplitude relative to the other. The exact zero condition will therefore combine Shannon's equality of probabilities with a phase opposition of their square roots.

\section{Entropy deficit as a balance cost}
Define the entropy deficit
\begin{equation}
D_S(\sigma)=\ln2-\Hbin(\sigma)\geq0.
\end{equation}
Near $\sigma=1/2$, write $\sigma=1/2+\delta$. A Taylor expansion gives
\begin{equation}
\Hbin\!\left(\frac12+\delta\right)
=\ln2-2\delta^2-\frac43\delta^4+O(\delta^6),
\end{equation}
so
\begin{equation}
D_S(\sigma)=2\delta^2+\frac43\delta^4+O(\delta^6).
\end{equation}
The leading penalty for leaving the balance axis is quadratic. Any operator model that realizes $D_S$ as a positive energy or action would have a unique minimum at the critical line. This observation will later motivate a candidate confinement term, while remaining distinct from a proof that zeta zeros minimize it.

\begin{figure}[ht]
\centering
\includegraphics[width=0.76\textwidth]{binary_entropy.pdf}
\caption{Binary Shannon entropy in nats. The unique balance point is $\sigma=1/2$, where the entropy equals $\ln2$.}
\label{fig:entropy}
\end{figure}
